\begin{statementbox}
    \textbf{\textcolor{CStatement}{这个结论解决什么问题？}}
    遇到“分子、分母都是一次式”的函数，题目常问：\textbf{函数值到底能取哪些数，有没有一个数永远取不到？}
    如果每次都代入许多 $x$ 去猜，既慢又不完整。把分子拆成“分母的若干倍 $+$ 一个数”，就能很快找到答案。

    \medskip
    \begin{center}\begin{tikzpicture}[font=\small,>=stealth,line join=round]
\draw[gray!45,rounded corners=3pt] (0,0) rectangle (14.2,7.2);
\node at (7.1,6.6) {分离常数后，看清两个永远取不到的坐标};
\begin{scope}[shift={(3.8,2.8)},scale=.58]
\draw[->] (-3.2,0)--(4.4,0) node[right]{$x$};\draw[->] (0,-2.1)--(0,5.4) node[above]{$y$};
\draw[dashed,magenta] (1,-2)--(1,5.1);\draw[dashed,magenta] (-3.1,2)--(4.2,2);
\draw[teal!80!black,thick,domain=-3:.22,samples=65] plot(\x,{2+3/(\x-1)});
\draw[teal!80!black,thick,domain=2:4.1,samples=65] plot(\x,{2+3/(\x-1)});
\draw (1,2) circle(2pt);
\node[below] at (1,0){$1$};\node[above left] at (0,2){$2$};

\end{scope}
\node[align=left,anchor=west] at (8.1,3.9){$x\ne1$，$y\ne2$\\[10pt] 中心 $(1,2)$\\[10pt] $(x-1)(y-2)=3$};
\node[text width=13.2cm,align=center] at (7.1,.65) {示例 $y=(2x+1)/(x-1)=2+3/(x-1)$；虚线为两条渐近线。};
\end{tikzpicture}
\end{center}
\end{statementbox}
\clearpage
\begin{statementbox}
\textbf{\textcolor{CStatement}{适用条件}}
    \begin{description}[style=nextline, leftmargin=2.8em, labelsep=0.5em, itemsep=0.35em]
        \item[\textbf{条件 1（函数的形式）：}]
            \(y=\dfrac{ax+b}{cx+d}\)，其中 $a,b,c,d$ 都是实数，$c\ne0$，$ad\ne bc$。
        \item[\textbf{条件 2（$x$ 的范围）：}]
            题目没有额外限制 $x$；只要求分母不为 $0$，即 $x\ne-\dfrac dc$。
    \end{description}

    \textbf{\textcolor{CStatement}{三个结论}}
    \begin{description}[style=nextline, leftmargin=2.8em, labelsep=0.5em, itemsep=0.45em]
        \item[\textbf{结论一（拆开分式）：}]
            \[
                \boxed{y=\frac ac+\frac{bc-ad}{c(cx+d)}}.
            \]
            右边第一项是固定的数，第二项永远不为 $0$。
        \item[\textbf{结论二（求值域）：}]
            \[
                \boxed{y\ne\frac ac}\qquad
                \text{即值域为 }\left\{y\in\mathbb R\mid y\ne\frac ac\right\}.
            \]
            这里不只是说“取不到 $a/c$”，而是说\textbf{除它以外，其余实数都能取到}。
        \item[\textbf{结论三（看图像）：}]
            图像围绕点
            \[
                \boxed{\left(-\frac dc,\frac ac\right)}
            \]
            成中心对称；直线 $x=-d/c$ 上没有图像上的点，函数值也不会等于 $a/c$。
    \end{description}

    \textbf{\textcolor{CStatement}{不能直接套用的情况}}
    若 $ad=bc$，函数实际上是常数函数，值域只有 $\{a/c\}$；若题目再规定 $x>0$、$x\ge2$ 等，值域还要按新的 $x$ 范围重新判断。
\end{statementbox}
